\section{Rounding the occupations to integer counts}
\label{sec:rounding}

The tensor maps of Section~\ref{sec:maps} apply each split and each tripling
a whole number of times, whereas the occupations \(\lambda_{r,x}\) and
\(\mu_r\) of Proposition~\ref{prop:occupations} are rational numbers whose
denominators may be large. For every integer rounding scale \(K\ge1\), this
section rounds \(K\) times the expected transition counts of each row to
integers that keep every state balanced exactly, as the row
identities~\eqref{eq:row-flow} do. Then
every tensor factor that one relation produces is consumed by another, apart
from \(K\) copies of \(S^+_{a,a}\) that are only consumed and \(K\) copies of
\(\MM_2\) that are only produced (Proposition~\ref{prop:rounded-ledger}). The
price, at every scale, is a rounding loss of at most \(4J_a\), where
\(J_a=O(a^3)\) is the number of transition coordinates
(Proposition~\ref{prop:entropy-loss}). Clearing the denominators of the exact
certificate of Appendix~\ref{sec:exact-certificate} would instead require a
scale that grows at least exponentially.

\subsection{Rounding a circulation}
\label{subsec:circulation-rounding}

A \emph{circulation} on a finite directed multigraph assigns a nonnegative
value to each edge so that at every vertex the values of the entering edges
and of the leaving edges have the same sum; a self-loop both enters and
leaves its vertex. The following lemma is classical: it follows from
\citet[Theorem~2]{HK56}, since the incidence matrix of a directed graph is
totally unimodular, and rounding through a return edge is described by
\citet[Sections~2.2--3.1]{KangPayor15}. We include a short proof because it
fixes the counts used in the tensor maps and allows self-loops and opposing
edges.

\begin{lemma}[Rounding a circulation]\label{lem:circulation-rounding}
Let \(\varphi\) be a circulation with rational values on a finite directed
multigraph, which may have self-loops and opposing edges. There is an
integer circulation \(n\) on the same multigraph with
\[
 \lfloor\varphi_e\rfloor\le n_e\le\lceil\varphi_e\rceil
 \quad\text{for every edge }e,
\]
and it can be computed with exact rational and integer arithmetic. In
particular, \(|n_e-\varphi_e|<1\) for every edge, and
\(n_e=\varphi_e\) whenever \(\varphi_e\) is an integer.
\end{lemma}

\begin{proof}
Put \(\eta_e=\varphi_e-\lfloor\varphi_e\rfloor\in[0,1)\). Since \(\varphi\)
is a circulation, the out-minus-in balance \(\beta_v\) of the \(\eta_e\) at
a vertex \(v\), to which self-loops contribute nothing, equals the
in-minus-out balance of the integer parts \(\lfloor\varphi_e\rfloor\), so
\(\beta_v\) is an integer. Form a network with a new source \(v_+\) and a
new sink \(v_-\): each nonintegral edge with distinct endpoints becomes its
own arc of capacity one, and we add an arc \(v_+\to v\) of capacity
\(\beta_v\) when \(\beta_v>0\) and an arc \(v\to v_-\) of capacity
\(-\beta_v\) when \(\beta_v<0\). The values \(\eta_e\), with the added arcs
saturated, form a flow of value
\(\beta_{\mathrm{tot}}=\sum_{v:\,\beta_v>0}\beta_v\), because at each
vertex \(v\) the arc from \(v_+\) or to \(v_-\) supplies or absorbs the
balance \(\beta_v\); since \(\sum_v\beta_v=0\), \(\beta_{\mathrm{tot}}\) is
also the total capacity of the arcs at \(v_+\), and of those at \(v_-\).
Since all capacities are integers, an augmenting-path algorithm started at the zero
flow raises the value by a positive integer at every step and stops after at
most \(\beta_{\mathrm{tot}}\) steps with an integral flow \(g\). When it
stops, the vertices reachable from \(v_+\) in the residual network form a
cut of capacity equal to the value of \(g\); this capacity is at least the
value \(\beta_{\mathrm{tot}}\) of the fractional flow, so \(g\) saturates
all arcs at \(v_+\) and \(v_-\). Set \(n_e=\lfloor\varphi_e\rfloor+g_e\),
where \(g_e\in\{0,1\}\) is the flow of \(g\) on the arc of \(e\), and
\(g_e=0\) for integral edges and self-loops. At every vertex the \(g_e\) have
out-minus-in balance \(\beta_v\) and the integer parts have balance
\(-\beta_v\), so \(n\) is an integer circulation with the required bounds.
Every step uses exact rational and integer arithmetic.
\end{proof}

\subsection{The rounded counts}
\label{subsec:rounded-counts}

Fix \(a\ge2\) and \(K\ge1\), and use the notation of
Section~\ref{sec:flows}; sums \(\sum_{r,x}\) run over \(2\le r\le a\) and
\(1\le x\le3r-1\). We round one graph for each row. The \emph{graph of
row \(r\)} has the states \(0,\ldots,3r\) as vertices and three kinds of
edges, each with a \emph{scaled value}. For every interior state \(x\) and
every child \(y\in\mathcal Y_{r,x}\) there is a \emph{transition edge}
\(x\to y\) with scaled value \(K\lambda_{r,x}q_{r,x,y}\); it records the
splits at \(x\) whose child is \(y\). The \emph{tripling edge} \(3r\to1\) has
scaled value \(K\mu_r\) and records the triplings of row \(r\). The
\emph{return edge} \(0\to\xin_r\) has scaled value \(K\); it returns the
\(K\) units absorbed at state \(0\) to the injection state. These scaled
values form a circulation: an edge from \(x\) to \(y\) with value \(\varphi\)
contributes \(\varphi(\mathbf e_x-\mathbf e_y)\) to the out-minus-in
balances, so, since \(\sum_yq_{r,x,y}=1\), the transition and tripling edges
together contribute \(K\) times the left side of~\eqref{eq:row-flow}, which
is \(K(\mathbf e_{\xin_r}-\mathbf e_0)\), and the return edge contributes
\(K(\mathbf e_0-\mathbf e_{\xin_r})\). The graph has self-loops, at the
states \(x\) with \(h_{r,x}\) even, and opposing edges, such as
\(0\to\xin_r\) and \(\xin_r\to0\). A loop does not affect any balance,
but its count is part of the number of splits at \(x\) and of their type,
so it must be rounded too; for example, in row \(2\) with \(a=2\) the loop
at state \(2\) has scaled value \(14K/15\).

Apply Lemma~\ref{lem:circulation-rounding} to the graph of each row
separately. The return edge has the integer value \(K\), so it keeps this
value; we remove it after rounding. Let \(n_{r,x,y}\) be the rounded count
on the transition edge \(x\to y\), which is the number of splits at state
\(x\) of row \(r\) with child \(y\), and let \(z_r\) be the rounded count on
the tripling edge, which is the number of triplings in row \(r\). Put
\(w_{r,x}=\sum_{y\in\mathcal Y_{r,x}}n_{r,x,y}\), the number of splits at
\((r,x)\), and \(\bar w=\sum_{r,x}w_{r,x}\) and \(\bar z=\sum_{r=2}^az_r\).
Together with the return edge, the rounded counts form a circulation, so
they satisfy the rounded row identities
\begin{equation}\label{eq:rounded-row-flow}
 \sum_{x=1}^{3r-1}
 \Bigl(w_{r,x}\mathbf e_x
       -\sum_{y\in\mathcal Y_{r,x}}n_{r,x,y}\mathbf e_y\Bigr)
 +z_r(\mathbf e_{3r}-\mathbf e_1)
 =K(\mathbf e_{\xin_r}-\mathbf e_0)
 \qquad(2\le r\le a).
\end{equation}
The lemma gives the first two of the strict error bounds
\begin{equation}\label{eq:rounded-edge-errors}
 |n_{r,x,y}-K\lambda_{r,x}q_{r,x,y}|<1,
 \qquad |z_r-K\mu_r|<1,
 \qquad |w_{r,x}-K\lambda_{r,x}|<d_{r,x},
\end{equation}
and the third follows by adding the first over the \(d_{r,x}\) children of
\(x\), since \(\sum_yq_{r,x,y}=1\).

An edge with scaled value zero keeps the count zero, so \(w_{r,x}=0\)
whenever \(\lambda_{r,x}=0\). The split at \((r,x)\) is \emph{active} if
\(w_{r,x}>0\); only active splits are used in the tensor maps. The count
\(z_r\) may be any nonnegative integer, since Lemma~\ref{lem:tripling}
imposes no divisibility condition. Every rule for choosing augmenting paths
gives counts with these properties; Theorem~\ref{thm:effective} fixes one.

\subsection{The ledger}
\label{subsec:ledger}

To see that the relations of all rows fit together, we record in one
formal sum, the \emph{ledger}, every tensor factor that a relation consumes,
with a plus sign, and every factor that it produces, with a minus sign. It
shows that all intermediate factors cancel, and it determines the catalyst:
the factors that are both consumed and produced. Recall the sign
\(\sigma_r=(-1)^{a-r}\) of row \(r\) (Section~\ref{sec:flows}), and write \(b_r(x)=r-1+x\) for the column of
state \(x\); thus state \(x\) of row \(r\) stands for the tensor
\(S^{\sigma_r}_{r,b_r(x)}\). We record tensor factors by formal symbols
\([\sigma;r,b]\), one for each tensor \(S^\sigma_{r,b}\), subject to the
identifications
\begin{equation}\label{eq:symbol-identifications}
 [\sigma;r,b]=[-\sigma;b,r],\qquad
 [+;r,r]=[-;r,r],\qquad
 [\sigma;1,b]=[\MM_b],
\end{equation}
which follow the isomorphisms \(S^\sigma_{r,b}\cong S^{-\sigma}_{b,r}\),
\(S^+_{r,r}\cong S^-_{r,r}\) and \(S^\sigma_{1,b}\cong\MM_b\)
of~\eqref{eq:cyclic-identities}; the second is the case \(b=r\) of the
first. We work in the free abelian group generated by the resulting
equivalence classes, which we still call symbols. The classes \([\MM_b]\)
are the \emph{boundary symbols}.

In row \(r\), a split at state \(x\) consumes the symbol
\([\sigma_r;r,b_r(x)]\) and produces \([\sigma_r;r,b_r(y)]\) for its child
\(y\) (Lemma~\ref{lem:multiway}); a tripling consumes
\([\sigma_r;r,4r-1]\) and produces \([\sigma_r;r,r]\)
(Lemma~\ref{lem:tripling} and~\eqref{eq:cyclic-tripling-balance}). The
\emph{source multiplicity} \(\operatorname{src}[\sigma;r,b]\) and the
\emph{output multiplicity} \(\operatorname{out}[\sigma;r,b]\) of a symbol
are the numbers of times it is consumed and produced: the \(w_{r,x}\) splits
at \((r,x)\) add \(w_{r,x}\) to \(\operatorname{src}[\sigma_r;r,b_r(x)]\),
and those among them with child \(y\) add \(n_{r,x,y}\) to
\(\operatorname{out}[\sigma_r;r,b_r(y)]\); the \(z_r\) triplings of row
\(r\) add \(z_r\) to \(\operatorname{src}[\sigma_r;r,4r-1]\) and to
\(\operatorname{out}[\sigma_r;r,r]\). Summed over all symbols, both
multiplicities give \(\bar w+\bar z\). The \emph{catalyst multiplicities}
and the \emph{catalyst} are
\begin{equation}\label{eq:rounded-catalyst}
 c_{\sigma,r,b}=\min\bigl(\operatorname{src}[\sigma;r,b],
                          \operatorname{out}[\sigma;r,b]\bigr),
 \qquad
 \Cat=\bigotimes_{\sigma,r,b}(S^\sigma_{r,b})^{\otimes c_{\sigma,r,b}},
 \qquad \Cat_N=\Cat^{\otimes N}.
\end{equation}
Here the product, like every sum over \((\sigma,r,b)\) below, runs over the
symbols, with one representative \((\sigma,r,b)\) chosen in each class;
by~\eqref{eq:cyclic-identities}, another choice changes \(\Cat\) only by an
isomorphism.

\begin{proposition}[Rounded ledger]\label{prop:rounded-ledger}
The rounded counts satisfy
\begin{align}
 &\sum_{r=2}^a\sum_{x=1}^{3r-1}
 \Bigl(w_{r,x}[\sigma_r;r,b_r(x)]
       -\sum_{y\in\mathcal Y_{r,x}}n_{r,x,y}[\sigma_r;r,b_r(y)]\Bigr)
 \notag\\
 &\qquad+\sum_{r=2}^az_r
   \bigl([\sigma_r;r,4r-1]-[\sigma_r;r,r]\bigr)
 =K[+;a,a]-K[\MM_2].\label{eq:rounded-ledger}
\end{align}
Consequently, the source and output multiplicities agree for every symbol
except two: the source multiplicity of \([+;a,a]\) exceeds its output
multiplicity by \(K\), and \([\MM_2]\), the only boundary symbol that
occurs, is never consumed and is produced \(K\) times. Hence its catalyst
multiplicity is zero, the catalyst \(\Cat\) is a tensor product of factors
\(S^\sigma_{r,b}\) with \(r,b\ge2\), and
\begin{equation}\label{eq:catalyst-size}
 \sum_{\sigma,r,b}c_{\sigma,r,b}=\bar w+\bar z-K.
\end{equation}
In tensor terms, the relations together consume
\(\Cat\otimes(S^+_{a,a})^{\otimes K}\) and produce
\(\Cat\otimes(S^{\sigma_2}_{2,1})^{\otimes K}\cong\Cat\otimes\MM_{2^K}\).
\end{proposition}

\begin{proof}
\emph{Step 1: the ledger identity.} In row \(r\), map each basis vector
\(\mathbf e_x\) to the symbol \([\sigma_r;r,b_r(x)]\) and extend linearly.
Since \(b_r(3r)=4r-1\) and \(b_r(1)=r\), this map sends the left side
of~\eqref{eq:rounded-row-flow} to the terms of row \(r\) on the left side
of~\eqref{eq:rounded-ledger}, and it sends the right side to
\(K([\sigma_r;r,b_r(\xin_r)]-[\sigma_r;r,r-1])\). In the final row,
\(\xin_a=1\) and \(\sigma_a=+\), so this is \(K([+;a,a]-[+;a,a-1])\). A
lower row has \(\xin_r=2\) and gives
\(K([\sigma_r;r,r+1]-[\sigma_r;r,r-1])\). Because
\(\sigma_{r+1}=-\sigma_r\), the first identification
in~\eqref{eq:symbol-identifications} gives
\([\sigma_r;r,r+1]=[\sigma_{r+1};r+1,r]\); thus, for \(2\le r<a\), the
positive term of row \(r\) cancels the negative term of row \(r+1\). Apart
from naming \([+;a,a]=[-;a,a]\), the proof uses the signs only through this
alternation (Section~\ref{subsec:orientations}). Summing over all rows
leaves \(K[+;a,a]-K[\sigma_2;2,1]\), and
\([\sigma_2;2,1]=[-\sigma_2;1,2]=[\MM_2]\). This
proves~\eqref{eq:rounded-ledger}.

\emph{Step 2: the multiplicities.} By their definition, the left side
of~\eqref{eq:rounded-ledger} equals
\(\sum_{\sigma,r,b}(\operatorname{src}[\sigma;r,b]
 -\operatorname{out}[\sigma;r,b])\,[\sigma;r,b]\). Since \(a\ge2\), the
symbols \([+;a,a]\) and \([\MM_2]\) are distinct, so comparing coefficients
gives equal multiplicities for every other symbol, and differences \(K\) for
\([+;a,a]\) and \(-K\) for \([\MM_2]\). A consumed symbol is
\([\sigma_r;r,b]\) with \(r\ge2\) and \(b\ge r\) (either \(b=b_r(x)\) for an
interior state \(x\), or \(b=4r-1\)); it and every symbol identified with
it have both indices at least \(2\). Hence
\(\operatorname{src}[\MM_2]=0\) and \(\operatorname{out}[\MM_2]=K\). A
produced symbol is \([\sigma_r;r,b_r(y)]\) or \([\sigma_r;r,r]\), with
\(r\ge2\) and \(b_r(y)\ge r-1\), so an index \(1\) occurs only for \(r=2\)
and \(y=0\), which gives \([\sigma_2;2,1]=[\MM_2]\).

\emph{Step 3: the catalyst.} By Step~2, the catalyst multiplicity equals
the source multiplicity for every symbol except \([+;a,a]\), where it is
\(\operatorname{out}[+;a,a]=\operatorname{src}[+;a,a]-K\), and it equals
the output multiplicity for every symbol except \([\MM_2]\), where it is
\(\min(0,K)=0\). Hence every symbol with positive catalyst multiplicity is
consumed, so its indices are at least \(2\); summing over all symbols gives
\(\sum_{\sigma,r,b}c_{\sigma,r,b}=\bar w+\bar z-K\); and the consumed
symbols, counted with multiplicity, are those of \(\Cat\) and \(K\) copies
of \([+;a,a]\), while the produced ones are those of \(\Cat\) and \(K\)
copies of \([\MM_2]\). The tensor statement follows
by~\eqref{eq:cyclic-identities} and \(\MM_2^{\otimes K}\cong\MM_{2^K}\).
\end{proof}

For \(2\le r<a\), the symbol \([\sigma_r;r,r+1]=[\sigma_{r+1};r+1,r]\) of
Step~1 belongs to state \(2\) of row \(r\) and to state \(0\) of row
\(r+1\); a comparison of
the column ranges \([r-1,4r-1]\) shows that different rows share no other
symbol. The auxiliary matrix factors are not ledger symbols;
Section~\ref{sec:extraction} accounts for them separately.

\subsection{The rounding loss}
\label{subsec:entropy-loss}

Rounding changes the frequencies with which each split produces its
children, and so it changes the gain of the construction, which we measure
by the rounded potential \(G\) below: \(NG\) approximates the logarithm
\(\Delta_N\) of the ratio of the one-step target side to the one-step source
side of Section~\ref{sec:extraction}, by~\eqref{eq:type-error}. The term
\(K\log2\) of \(G\) comes from the output \(\MM_{2^{NK}}\); the term of a
split at \((r,x)\) compares the roughly \(\exp(Nw_{r,x}H(\hat q_{r,x}))\)
sectors of its \(Nw_{r,x}\) uses with their auxiliary side
\(d_{r,x}^{Nw_{r,x}}\); and \(z_r\log3\) counts the roughly \(3^{Nz_r}\)
sectors of the triplings of row \(r\).

For probability vectors \(\hat q\) and \(q\) on
the same finite set, with \(q\) positive, let
\(H(\hat q)=-\sum_y\hat q_y\log\hat q_y\), with \(0\log0=0\), be the
entropy, and let
\(D(\hat q\Vert q)=\sum_{y:\,\hat q_y>0}\hat q_y\log(\hat q_y/q_y)\) be the
relative entropy, which is nonnegative (Gibbs' inequality). For an active
split at \((r,x)\), let \(\hat q_{r,x}=(\hat q_{r,x,y})_y\), with
\(\hat q_{r,x,y}=n_{r,x,y}/w_{r,x}\), be its empirical distribution of
children, and let \(q_{r,x}=(q_{r,x,y})_y\). The \emph{rounded potential}
and the \emph{rounding loss} are
\begin{equation}\label{eq:rounded-gain}
 G=K\log2+\sum_{r,x:\,w_{r,x}>0}
   w_{r,x}\bigl(H(\hat q_{r,x})-\log d_{r,x}\bigr)+\bar z\log3,
 \qquad
 \mathcal L=\sum_{r,x:\,w_{r,x}>0}w_{r,x}D(\hat q_{r,x}\Vert q_{r,x}).
\end{equation}
Finally, let \(J_a\) be the number of transition coordinates, that is, the
number of transition edges in all row graphs, loops included:
\begin{equation}\label{eq:coordinate-count}
 \begin{split}
 J_a&=\sum_{r=2}^a\sum_{x=1}^{3r-1}d_{r,x}
     =\sum_{r=2}^a\Bigl(\Bigl\lfloor\frac{9r^2}{4}\Bigr\rfloor+3r-1\Bigr)\\
    &=\frac{12a^3+42a^2+12a-65+(-1)^a}{16}=O(a^3).
 \end{split}
\end{equation}
The second equality holds because
\(\sum_{x=1}^{3r-1}\min(x,3r-x)=\lfloor9r^2/4\rfloor\). The third follows by
summing over \(r\), using \(\lfloor9r^2/4\rfloor=9r^2/4\) for even \(r\) and
\(\lfloor9r^2/4\rfloor=(9r^2-1)/4\) for odd \(r\).

\begin{proposition}[Uniform entropy loss]\label{prop:entropy-loss}
For every rounding scale \(K\ge1\), the rounded potential and the rounding
loss of~\eqref{eq:rounded-gain} satisfy
\begin{equation}\label{eq:rounded-entropy}
 G=K\log F(a,a)-\mathcal L,
 \qquad 0\le\mathcal L\le4J_a.
\end{equation}
\end{proposition}

\begin{proof}
The identity comes from evaluating the potential \(\log F\) on the ledger.
The bound comes from comparing each empirical distribution with the
transition probabilities edge by edge, using that no child probability is
smaller than \(1/(3d_{r,x})\).

\emph{Step 1: a lower bound for the child probabilities.} Fix a row \(r\)
and an interior state \(x\), and write \(h=h_{r,x}\) and \(d=d_{r,x}\).
By~\eqref{eq:multiway-probabilities},
\(q_{r,x,y}=f_r(y)/(d\,f_r(x))\), where \(f_r(y)=\frac{3(r-1)}2+y\) is
increasing; so the smallest probability belongs to the smallest child
\(x-h\) and equals \(f_r(x-h)/(d\,f_r(x))\). If \(x\le3r/2\), then \(h=x\),
and the ratio \(f_r(0)/f_r(x)\) decreases with \(x\). If \(x\ge3r/2\), then
\(h=3r-x\), and the ratio
\(f_r(2x-3r)/f_r(x)=2-(f_r(0)+3r)/f_r(x)\) increases with \(x\). Both
expressions, read as functions of a real variable \(x\), take the common
value \(f_r(0)/f_r(3r/2)=(r-1)/(2r-1)\) at \(x=3r/2\), whether or not \(3r/2\)
is an integer. Therefore every child \(y\) of \(x\) satisfies
\begin{equation}\label{eq:child-probability-bound}
 q_{r,x,y}\ge\frac{f_r(0)}{d_{r,x}f_r(3r/2)}
 =\frac{r-1}{(2r-1)d_{r,x}}\ge\frac1{3d_{r,x}},
\end{equation}
where the last step uses \(r\ge2\). The bound is attained for \(r=2\) and
\(x=3\), where \(q_{2,3,0}=1/12\).

\emph{Step 2: the loss of one split.} Fix an active split at \((r,x)\). In
this step we omit the subscripts \(r,x\): thus \(\lambda=\lambda_{r,x}\),
\(w=w_{r,x}\) and \(d=d_{r,x}\), and \(n_y\), \(\hat q_y\), \(q_y\) stand
for \(n_{r,x,y}\), \(\hat q_{r,x,y}\), \(q_{r,x,y}\).
By~\eqref{eq:rounded-edge-errors}, \(n_y<K\lambda q_y+1\) for every child
\(y\), and \(K\lambda<w+d\); by~\eqref{eq:child-probability-bound},
\(1/q_y\le3d\). Hence every child with \(n_y>0\) satisfies
\[
 \frac{\hat q_y}{q_y}=\frac{n_y}{wq_y}
 <\frac{K\lambda}{w}+\frac1{wq_y}
 <1+\frac dw+\frac{3d}w=1+\frac{4d}w.
\]
Averaging the logarithm of this inequality with the weights \(\hat q_y\)
gives \(D(\hat q_{r,x}\Vert q_{r,x})\le\log(1+4d/w)\), and therefore
\[
 w\,D(\hat q_{r,x}\Vert q_{r,x})\le w\log\Bigl(1+\frac{4d}w\Bigr)\le4d.
\]
The argument divides only by the positive integer \(w\), never by the
occupation \(\lambda\), so the bound holds at every scale \(K\), even when
some occupations are tiny or zero. Summing over the active splits, and
using \(D\ge0\), gives \(0\le\mathcal L\le\sum_{r,x}4d_{r,x}=4J_a\).

\emph{Step 3: the identity.} By~\eqref{eq:affine-row},
\(F(r,b_r(y))=s_rf_r(y)\) for \(0\le y\le3r\). Hence the transition
probabilities~\eqref{eq:multiway-probabilities} are ratios of profile values,
\[
 q_{r,x,y}=\frac{F(r,b_r(y))}{d_{r,x}F(r,b_r(x))},
 \quad\text{that is,}\quad
 \log F(r,b_r(y))=\log F(r,b_r(x))+\log d_{r,x}+\log q_{r,x,y}.
\]
Multiply by \(n_{r,x,y}\) and sum over the children. For an active split,
\(\sum_yn_{r,x,y}=w_{r,x}\), and the definitions of \(H\) and \(D\) give
\(-\sum_{y:\,n_{r,x,y}>0}\hat q_{r,x,y}\log q_{r,x,y}
 =H(\hat q_{r,x})+D(\hat q_{r,x}\Vert q_{r,x})\). Hence
\begin{equation}\label{eq:split-potential}
 w_{r,x}\log F(r,b_r(x))-\sum_yn_{r,x,y}\log F(r,b_r(y))
 =w_{r,x}\bigl(H(\hat q_{r,x})-\log d_{r,x}\bigr)
  +w_{r,x}D(\hat q_{r,x}\Vert q_{r,x}).
\end{equation}
For an inactive split both sides are zero. For a tripling,
\(F(r,4r-1)=3F(r,r)\) (Step~1 of the proof of
Theorem~\ref{thm:least-profile}), so
\begin{equation}\label{eq:tripling-potential}
 z_r\bigl(\log F(r,4r-1)-\log F(r,r)\bigr)=z_r\log3.
\end{equation}
Now assign to each symbol \([\sigma;r,b]\) the value \(\log F(r,b)\). This
is well defined, because \(F\) is symmetric, the value does not depend on
\(\sigma\), and \(F(1,b)=b\), so that the boundary symbol
\([\MM_b]=[\sigma;1,b]\) receives \(\log b\). Extend the assignment
linearly and apply it to~\eqref{eq:rounded-ledger}.
By~\eqref{eq:split-potential} and~\eqref{eq:tripling-potential}, the left
side becomes
\[
 \sum_{r,x:\,w_{r,x}>0}w_{r,x}\bigl(H(\hat q_{r,x})-\log d_{r,x}\bigr)
 +\mathcal L+\bar z\log3,
\]
while the right side becomes \(K\log F(a,a)-K\log2\). Adding \(K\log2\) to
both sides gives \(G=K\log F(a,a)-\mathcal L\).
\end{proof}

Thus \(w_{r,x}D(\hat q_{r,x}\Vert q_{r,x})\) is the price of producing
children with the frequencies \(\hat q_{r,x}\) instead of \(q_{r,x}\); it
vanishes if \(K\) clears the denominators of all scaled values. For a
distribution \(\hat p\) with entries in \(\frac1{w_{r,x}}\mathbb Z\) that
minimizes \(D(\hat p\Vert q_{r,x})\), the bound of \citet[Section~IV.D, Eq.~(42)]{BoechererGeiger16}
and~\eqref{eq:child-probability-bound} give \(3d_{r,x}\) in place of
\(4d_{r,x}\), and so does Step~2 when \(w_{r,x}\ge K\lambda_{r,x}\); the extra
\(d_{r,x}\) pays for the splits with \(w_{r,x}<K\lambda_{r,x}\).

\begin{remark}[Size of the loss in examples]\label{rem:observed-loss}
The bound \(4J_a\) is far from tight. A floating-point evaluation for the
rounding rule of Theorem~\ref{thm:effective}, with \(2\le a\le12\) and
\(K\in\{1,2,7,a^3,a^3\lceil\log_2a\rceil\}\), gives
\(\mathcal L/(4J_a)<0.029\) in every case. The largest ratio occurs at
\(a=2\) and \(K=1\), where \(\mathcal L=\log5\) and \(4J_2=56\). For
\(a=12\) and the scale \(K=a^3\lceil\log_2a\rceil=6912\) of
Theorem~\ref{thm:quantitative} (Section~\ref{sec:costs} uses a larger
\(K\)), \(\mathcal L\approx3.96\), against \(4J_{12}=6716\).
\end{remark}

\subsection{Bounds on the rounded counts}
\label{subsec:count-bounds}

Put \(W=\sum_{r,x}w_{r,x}\log d_{r,x}\); at type scale \(N\), the source
auxiliary matrix factor of all splits has side \(\exp(NW)\). Recall
from~\eqref{eq:occupation-bounds} that \(\sum_{r,x}\lambda_{r,x}\) and
\(\sum_{r,x}\lambda_{r,x}\log d_{r,x}\) are both at most \(\Gamma_a\), and
from Proposition~\ref{prop:occupations} that
\(\lambda_{r,x}\log d_{r,x}\le180\) for every split. These bounds
and~\eqref{eq:rounded-edge-errors} give
\begin{equation}\label{eq:rounded-count-bounds}
 \begin{split}
 W&\le K\Gamma_a+J_a\log(2a),\qquad \bar w\le K\Gamma_a+J_a,\\
 \bar z&\le(2K+1)(a-1),\\
 w_{r,x}\log d_{r,x}&\le180K+2a\log(2a).
 \end{split}
\end{equation}
Indeed, \(d_{r,x}\le\lfloor3a/2\rfloor+1\le2a\). Summing
\(w_{r,x}<K\lambda_{r,x}+d_{r,x}\) over all splits gives
\(\bar w<K\sum_{r,x}\lambda_{r,x}+J_a\). Multiplying the inequality \(w_{r,x}<K\lambda_{r,x}+d_{r,x}\) of a single
split by \(\log d_{r,x}\le\log(2a)\) gives
\(w_{r,x}\log d_{r,x}<K\lambda_{r,x}\log d_{r,x}+d_{r,x}\log(2a)\); with
\(\lambda_{r,x}\log d_{r,x}\le180\) this proves the last bound, and summing
over all splits proves the bound for \(W\). By~\eqref{eq:reset-weights},
\(\mu_r\le2\), so \(z_r<2K+1\) in each of the \(a-1\) rows. The bound for a
single split controls the largest sector count in Section~\ref{sec:rate};
the bound for \(W\) alone would be weaker by a factor of order \(a\).
